%%% SOURCE DOCUMENTS AND THEIR SIGNIFICANCE %%%
\thispagestyle{acknowledgements}

This National Platform draws upon a comprehensive body of completed work
spanning regulatory adaptations, clinical trial standards, patient-centered
documentation, national infrastructure designs, and supporting research. Each
source document listed below is referenced throughout this paper by its
designated name. The following overview defines the role and significance of
each document within the platform.

\subsection*{Regulatory Adaptation Documents}

\noindent\textbf{Adaption: ICH Harmonised Guideline} \cite{ich-adapt}.
A complete adaptation of the ICH E6(R3) Guideline for Good Clinical Practice
to Physical AI oncology clinical trials. This document establishes principles
of Physical AI clinical practice, investigator responsibilities, sponsor
responsibilities, and data governance requirements across four major sections
and three appendices. It is referenced throughout Sections~\ref{sec:ich-e6r3},
\ref{sec:site-establishment}, and \ref{sec:discussion} as the foundational
good clinical practice standard for Physical AI implementations.

\vspace{0.3cm}
\noindent\textbf{Physical AI Adaption: 21 CFR Part 50} \cite{cfr50-adapt}.
An adaptation of the federal regulation governing Protection of Human Subjects
to Physical AI clinical trial contexts. This document addresses informed
consent requirements specific to robotic and AI-assisted procedures, IRB review
standards for Physical AI investigations, pediatric protections, and a
comprehensive Physical AI system classification framework. It is referenced in
Sections~\ref{sec:cfr50}, \ref{sec:patient-journey},
and~\ref{sec:patient-instructions} as the human subjects protection standard.

\vspace{0.3cm}
\noindent\textbf{Physical AI Adaption: 21 CFR Part 312} \cite{cfr312-adapt}.
An adaptation of the Investigational New Drug Application regulation to Physical
AI trial contexts spanning ten subparts, including a new Subpart J dedicated
entirely to Physical AI systems. This document covers IND requirements, clinical
phases including a Phase 0 simulation validation stage, safety reporting for
Physical AI adverse events and cybersecurity incidents, and expanded access
provisions. It is referenced in Sections~\ref{sec:cfr312},
\ref{sec:site-establishment}, and~\ref{sec:implementation}.

\subsection*{Standards and Framework Documents}

\noindent\textbf{Physical AI Unification Standard Level} \cite{usl-standard}.
Defines the Unification Standard Level (USL) scoring methodology for evaluating
robot readiness across four dimensions: simulation framework switching, AI
integration, cross-robot sharing, and clinical trial collaboration. This
document provides quantitative scoring for nine robots across three categories
(cobots, surgical, humanoids) used in Physical AI oncology trials. It is
referenced throughout Sections~\ref{sec:psl-usl},
\ref{sec:patient-journey}, and~\ref{sec:discussion} as the robot readiness
evaluation framework.

\vspace{0.3cm}
\noindent\textbf{Clinical Trial Site Complete Documentation} \cite{site-docs}.
A comprehensive package of eleven regulatory documents required for establishing
the first Physical AI oncology clinical trial site in California. This includes
state legislation (SB 1042, AB 2847, SB 892), municipal code updates, Title 22
regulations, FDA compliance guides, building and premises codes, parking and
transportation standards, activation SOPs, and emergency preparedness plans. The
Physical AI Standard Level (PSL) framework with its three regulatory dimensions
is defined within this body of work. Referenced in Sections~\ref{sec:psl-usl},
\ref{sec:site-establishment}, and~\ref{sec:implementation}.

\subsection*{Patient-Centered Documents}

\noindent\textbf{A Cancer Patient's Journey, Physical AI}
\cite{patient-journey-paper}.
A comprehensive paper documenting a fully autonomous single-patient journey
through a regulated Physical AI oncology trial across ten stages, from
prescreening through closeout. This document includes treatment outcomes, FDA
cost-saving analysis, USL scoring integration, and thirty ASCII diagrams
illustrating the patient pipeline. It serves as primary evidence for Physical AI
feasibility and is referenced in Sections~\ref{sec:patient-journey},
\ref{sec:financial}, and~\ref{sec:discussion}.

\vspace{0.3cm}
\noindent\textbf{Patient Instructions: Physical AI Trials}
\cite{patient-instructions-paper}.
A ten-page patient-facing document providing clear, accessible instructions for
participants in Physical AI oncology trials. It covers interactions with ten
robot types including surgical robots, cobots, radiotherapy positioning robots,
social companion robots, humanoids, and rehabilitation exoskeletons. Referenced
in Sections~\ref{sec:patient-instructions} and~\ref{sec:patient-journey} as the
patient communication standard.

\subsection*{National Infrastructure Documents}

\noindent\textbf{Physical AI National MCP Servers} \cite{national-mcp-paper}.
Describes the design, implementation, and testing of five national Model Context
Protocol (MCP) servers for Physical AI oncology clinical trial systems. This
document covers the complete server architecture, tool inventories, conformance
levels, safety modules, integration adapters, and national deployment topology.
Referenced in Sections~\ref{sec:national-mcp},
\ref{sec:federated-learning}, and~\ref{sec:implementation}.

\vspace{0.3cm}
\noindent\textbf{Physical AI Federated Learning} \cite{fl-paper}.
Presents a federated learning framework for unifying Physical AI oncology trials
across multiple sites while preserving patient privacy. This document covers
architecture design, five foundational pillars, triple AI peer review,
code trust and safety, clinical analytics, and digital twin integration.
Referenced in Sections~\ref{sec:federated-learning}, \ref{sec:national-mcp},
and~\ref{sec:discussion}.

\subsection*{Supporting Research}

\noindent\textbf{Research A: U.S. Government Branch Operations.}
A structured inventory of constitutional provisions and federal statutes
governing the three branches of U.S. government, with detailed analysis of how
AI in clinical trials is regulated through existing legislative, executive, and
judicial frameworks. This research informs the governance structure adapted to
Physical AI oncology trials in Section~\ref{sec:gov-framework}.
\textit{Source files: national-platform/research\_a/01\_research\_a\_part1.txt
and 02\_research\_a\_part2.txt.}

\vspace{0.3cm}
\noindent\textbf{Research B: California and Federal Regulatory Analysis.}
A dual-jurisdiction analysis of California state law and federal law as they
apply to AI and robotics in oncology clinical trials. This research catalogs
regulatory requirements from both jurisdictions, identifies gaps and
ambiguities, and informs the regulatory landscape analysis in
Section~\ref{sec:regulatory-landscape}. It also serves as fact-checking
material for bill names and governance structures referenced in the Clinical
Trial Site Complete Documentation.
\textit{Source files: national-platform/research\_b/01\_research\_b\_part1.txt
and 02\_research\_b\_part2.txt.}

\subsection*{Code Repositories}

\noindent\textbf{Physical AI for Oncology Clinical Trials Repository}
\cite{main-repo}.
The primary repository containing all source code, simulation frameworks,
orchestration modules, digital twin implementations, regulatory compliance
tools, and documentation for the Physical AI oncology trials platform.

\vspace{0.3cm}
\noindent\textbf{National MCP Servers Repository} \cite{mcp-repo}.
The repository containing the five national MCP server implementations,
integration adapters, and deployment configurations for Physical AI oncology
clinical trial systems.

\vspace{0.3cm}
\noindent\textbf{PAI Oncology Trial FL Repository} \cite{fl-repo}.
The repository containing the federated machine learning framework, privacy
modules, and multi-site coordination tools for Physical AI oncology trials.
